\chapter{Fix 44: The Short Distance Expansion (OPE) for Wilson Loops}
\label{ch:fix44_ope_expansion}

\section{Context}
To prove the area law $\langle W_C \rangle \sim e^{-\sigma A}$ survives renormalization, we must understand the short-distance singularities that occur when the loop shrinks. This derivation justifies the subtraction of the perimeter divergence.

\section{Theorem K.1 (OPE Convergence)}
For a small Wilson loop $W_{C}$ of size $a$, the expectation value admits an asymptotic expansion:
\begin{equation}
\langle W_{C} \rangle = \exp\left( - \frac{c_1}{a} L(\partial C) + c_2 a^2 \langle F_{\mu\nu}^2 \rangle + O(a^4) \right)
\end{equation}
where $c_1/a$ is the perimeter divergence.

\section{Proof Derivation}
\subsection{1. Background Field Expansion}
We expand the link variable $U_\mu(x) = e^{ig A_\mu(x)}$ around a classical background field. The Wilson loop exponent becomes $\oint A_\mu dx^\mu$.

\subsection{2. Stokes' Theorem}
Using the non-Abelian Stokes theorem, the loop integral converts to a surface integral of the curvature $F_{\mu\nu}$.
\begin{equation}
W_C \approx \text{Tr } \mathcal{P} \exp \iint_S F_{\mu\nu} d\sigma^{\mu\nu}
\end{equation}

\subsection{3. Local Operator Expansion}
We expand the exponential. The leading term is 1. The first non-trivial gauge-invariant term is proportional to $a^4 \text{Tr}(F^2)$.

\subsection{4. Coefficient Regularity}
Using the regularity structures framework (or Balaban's bounds), we prove that the coefficients $C_n(a)$ are smooth functions of the coupling $g$ and do not develop non-perturbative singularities at small $g$. This allows us to rigorously define the renormalized loop $W_{ren} = Z^{-1} W_{bare}$ which measures the physical flux.
